\chapter{Methodology}

This chapter presents the detailed methodology of the PISL framework as proposed by Tao~\textit{et~al.} \cite{tao2024unsupervised}. The overall architecture is illustrated in Figure~\ref{fig:overview}.

\begin{figure}[H]
    \centering
    \includegraphics[width=\textwidth]{overview.png}
    \caption{Overview of the PISL framework \cite{tao2024unsupervised}. The framework consists of a ResNet-50 backbone, a Spatial Diffusion Model (SDM) that generates transformation parameters through denoising, and a sampler that extracts semantically meaningful patches. Training is guided by five losses: $\mathcal{L}_{ces}$, $\mathcal{L}_{tri}$, $\mathcal{L}_{cam}$, $\mathcal{L}_{scd}$, and $\mathcal{L}_{psc}$.}
    \label{fig:overview}
\end{figure}

\section{Overall Framework}

Given an unlabelled dataset $\mathcal{X} = \{x_i\}_{i=1}^{N}$ with camera annotations $\{c_i\}_{i=1}^{N}$ but without identity labels, the PISL framework iteratively performs four steps:

\begin{enumerate}
    \item \textbf{Feature Extraction:} Extract global features $\mathbf{f}_g \in \mathbb{R}^{D}$ and part features $\mathbf{f}_p \in \mathbb{R}^{D \times P}$ using the backbone and Spatial Diffusion Model, where $D$ is the feature dimension and $P$ is the number of parts.
    \item \textbf{Pseudo Label Generation:} Cluster global features using DBSCAN \cite{ester1996dbscan} with Jaccard distance to assign pseudo identity labels.
    \item \textbf{Semantic Consistency Scoring:} Compute consistency scores between global and part features using Maximum Mean Discrepancy (MMD) \cite{gretton2012mmd}.
    \item \textbf{Model Training:} Update the model using a combination of classification, metric learning, camera-aware contrastive, and diffusion-based losses.
\end{enumerate}

\section{Backbone and Feature Extraction}

The backbone network is a ResNet-50 \cite{he2016deep} pretrained on ImageNet \cite{deng2009imagenet}. Given an input image $x$, the backbone extracts a feature map $\mathbf{F} \in \mathbb{R}^{2048 \times H \times W}$ from the last convolutional layer. A global average pooling operation produces the global feature vector:
\begin{equation}
    \mathbf{f}_g = \text{GAP}(\mathbf{F}) \in \mathbb{R}^{2048}
\end{equation}
which is subsequently normalized via batch normalization and L2 normalization for distance computation.

\section{Spatial Diffusion Model (SDM)}

The Spatial Diffusion Model is the central innovation of PISL. Instead of generating patches through a learned static spatial transformer, PISL introduces a \textit{diffusion process} over the spatial transformer parameters $\boldsymbol{\theta}$.

\subsection{Localization Network}

A localization network processes the backbone feature map $\mathbf{F}$ to produce initial transformation parameters $\boldsymbol{\theta}_0 \in \mathbb{R}^{P \times 2 \times 3}$, where $P$ is the number of parts (set to 3 in practice). The localization network consists of:
\begin{itemize}
    \item A convolutional layer: Conv2d($2048 \rightarrow 4096$, kernel $3 \times 3$)
    \item Batch normalization and ReLU activation
    \item Max pooling (kernel $3 \times 3$, stride 2)
    \item Adaptive average pooling to $1 \times 1$
    \item Two fully connected layers: FC($4096 \rightarrow 512$) $\rightarrow$ FC($512 \rightarrow 3 \times 2 \times 3$)
\end{itemize}

Each $\boldsymbol{\theta}_0^{(i)}$ defines an affine transformation for the $i$-th part:
\begin{equation}
    \boldsymbol{\theta}_0^{(i)} = \begin{pmatrix} s_x & 0 & t_x \\ 0 & s_y & t_y \end{pmatrix}
\end{equation}
where $s_x, s_y$ control the scale and $t_x, t_y$ control the translation. The parameters are initialized to produce three vertically stacked horizontal strips corresponding to head, torso, and legs.

\subsection{Forward Diffusion Process}

The forward process gradually adds Gaussian noise to $\boldsymbol{\theta}_0$ over $T = 1000$ timesteps:
\begin{equation}
    q(\boldsymbol{\theta}_t | \boldsymbol{\theta}_{t-1}) = \mathcal{N}\left(\boldsymbol{\theta}_t;\ \sqrt{1-\beta_t}\,\boldsymbol{\theta}_{t-1},\ \beta_t\,\mathbf{I}\right)
\end{equation}

Using the reparameterization trick with $\alpha_t = 1 - \beta_t$ and $\bar{\alpha}_t = \prod_{s=1}^{t}\alpha_s$:
\begin{equation}
    \boldsymbol{\theta}_t = \sqrt{\bar{\alpha}_t}\,\boldsymbol{\theta}_0 + \sqrt{1 - \bar{\alpha}_t}\,\boldsymbol{\epsilon}, \quad \boldsymbol{\epsilon} \sim \mathcal{N}(\mathbf{0}, \mathbf{I})
\end{equation}

A \textit{cosine} $\beta$-schedule \cite{ho2020denoising} is employed with offset $s = 0.008$:
\begin{equation}
    \bar{\alpha}_t = \frac{f(t)}{f(0)}, \quad f(t) = \cos\left(\frac{t/T + s}{1 + s} \cdot \frac{\pi}{2}\right)^2
\end{equation}

\subsection{Reverse Denoising Process}

A denoising network $\boldsymbol{\Phi}$ predicts the noise $\boldsymbol{\epsilon}$ at each timestep:
\begin{equation}
    \hat{\boldsymbol{\theta}}_0 = \frac{1}{\sqrt{\bar{\alpha}_t}} \left( \boldsymbol{\theta}_t - \sqrt{1 - \bar{\alpha}_t} \cdot \boldsymbol{\Phi}(\boldsymbol{\theta}_t, t) \right)
\end{equation}

\subsection{Patch Sampling via Spatial Transformer}

The denoised parameters $\hat{\boldsymbol{\theta}}_0$ construct a sampling grid via the STN \cite{jaderberg2015spatial}:
\begin{equation}
    \mathbf{G}_i = \hat{\boldsymbol{\theta}}_0^{(i)} \cdot \mathbf{G}_{\text{base}}
\end{equation}
Bilinear sampling extracts part feature maps from the backbone feature map $\mathbf{F}$, which are processed through global average pooling and batch normalization to produce part feature vectors $\mathbf{f}_p^{(i)} \in \mathbb{R}^{D}$.

\section{Semantic Controlled Diffusion (SCD) Loss}

The SCD loss guides the denoising process toward semantically meaningful transformations and consists of three components.

\subsection{Noise Prediction Loss}
\begin{equation}
    \mathcal{L}_{\text{noise}} = \|\boldsymbol{\Phi}(\boldsymbol{\theta}_t, t) - \boldsymbol{\epsilon}\|_2^2
\end{equation}

\subsection{Theta Consistency Loss}
Combines L2 distance and cosine similarity to keep denoised parameters close to the localization network's output:
\begin{equation}
    \mathcal{L}_{\text{consist}} = \|\hat{\boldsymbol{\theta}}_0 - \boldsymbol{\theta}_0\|_2^2 + 0.1 \cdot \left(1 - \cos(\hat{\boldsymbol{\theta}}_0, \boldsymbol{\theta}_0)\right)
\end{equation}

\subsection{Part Contrastive Loss}
Encourages discriminative part features with learned part weights:
\begin{equation}
    \mathcal{L}_{\text{part}} = -\sum_{i=1}^{P} w_i \cdot \frac{1}{|\mathcal{P}_i|} \sum_{j \in \mathcal{P}_i} \log \frac{\exp(\mathbf{f}_p^{(i)} \cdot \mathbf{f}_{p,j}^{(i)} / \tau)}{\sum_k \exp(\mathbf{f}_p^{(i)} \cdot \mathbf{f}_{p,k}^{(i)} / \tau)}
\end{equation}
where $w_i = \exp(-\mathcal{L}_{\text{part}}^{(i)}) / \sum_j \exp(-\mathcal{L}_{\text{part}}^{(j)})$ and $\tau = 0.07$.

The combined SCD loss is:
\begin{equation}
    \mathcal{L}_{scd} = \mathcal{L}_{\text{noise}} + \lambda_c \cdot \mathcal{L}_{\text{consist}} + \lambda_p \cdot \mathcal{L}_{\text{part}}
\end{equation}
where $\lambda_c = 0.5$ and $\lambda_p = 0.3$.

\section{Patch Semantic Consistency (PSC) Loss}

\subsection{Consistency Measurement via MMD}
For each sample $j$ and part $i$, the $k$-nearest neighbour sets are computed in both global ($\mathcal{G}_j$) and part ($\mathcal{P}_j^{(i)}$) feature spaces:
\begin{equation}
    s_j^{(i)} = \exp\left(-\text{MMD}^2(\mathcal{G}_j, \mathcal{P}_j^{(i)})\right)
\end{equation}
A high consistency score indicates agreement between global and part neighbourhood structures.

\subsection{PSC Label Refinement (PSC-LR)}
Uses weighted part predictions to refine global pseudo labels:
\begin{equation}
    \tilde{y} = \lambda_{\text{ref}} \cdot y + (1 - \lambda_{\text{ref}}) \cdot \sum_{i=1}^{P} w_i \cdot p_i
\end{equation}
where $y$ is the one-hot pseudo label, $p_i$ is the softmax prediction from part $i$, and $\lambda_{\text{ref}} = 0.5$.

\subsection{PSC Label Smoothing (PSC-LS)}
Applies consistency-aware label smoothing:
\begin{equation}
    \mathcal{L}_{\text{psc-ls}}^{(i)} = s_j^{(i)} \cdot \mathcal{L}_{ce}(\mathbf{f}_p^{(i)}, y) + (1 - s_j^{(i)}) \cdot D_{KL}(p_i \,\|\, u)
\end{equation}
where $u$ is the uniform distribution. High consistency emphasizes cross-entropy; low consistency regularizes toward uniform predictions.

\section{Camera-Aware Contrastive Learning}

Camera-level proxies are constructed for each (pseudo identity, camera) pair \cite{wang2021cap}:
\begin{equation}
    \mathcal{L}_{cam} = -\sum_{i=1}^{B} \frac{1}{|\mathcal{P}_i^+|} \sum_{j \in \mathcal{P}_i^+} \log \frac{\exp(s_{ij}/\tau)}{\sum_{k \in \mathcal{P}_i^+ \cup \mathcal{N}_i^-} \exp(s_{ik}/\tau)}
\end{equation}
where $\mathcal{P}_i^+$ contains positive proxies of the same identity but different cameras, $\mathcal{N}_i^-$ contains the top-50 hard negative proxies, and $\tau = 0.07$.

\section{Total Training Objective}

\begin{equation}
    \mathcal{L} = \mathcal{L}_{ces} + \mathcal{L}_{tri} + W_{cam} \cdot \mathcal{L}_{cam} + W_{diff} \cdot \mathcal{L}_{scd} + \mathcal{L}_{psc}
\end{equation}
where $\mathcal{L}_{ces}$ is cross-entropy with label smoothing \cite{szegedy2016rethinking}, $\mathcal{L}_{tri}$ is the soft triplet loss \cite{hermans2017triplet}, $W_{cam} = 0.5$, and $W_{diff} = 0.1$.
